\section{Related work}

The closest supervised comparison is the unrestricted discrete model of \citet[model (3), Theorems 1--2]{ZengBalakrishnanHanKennedy2026Discrete}. For a fixed overlap floor, their conventional estimators have squared-error scale \(n^{-1}+d^2/n^2\), while their converse has scale \(n^{-1}+d^2/(n^2\log^2 n)\) in its stated many-cell range, with floor-dependent constants. The public CausalSmith working paper \citet[Theorem \texttt{thm:sharp-minimax-fixed-interior}]{InternalDiscreteATEPolynomialUpperMatch} closes that gap: for each fixed \(0<\epsilon<1/2\), there are \(a_\epsilon,\rho_\epsilon,C_\epsilon>0\) and \(N_\epsilon<\infty\) such that, whenever \(n\ge N_\epsilon\) and \(d\le\rho_\epsilon n\log n\), the supervised minimax risk is bounded above and below by floor-dependent multiples of
\[
n^{-1}+\frac{d^2}{n^2\log^2 n}.
\]
Our comparison in \cref{obj:thm:uniform-annotation-frontier} treats the complete-record budget, auxiliary-record budget, alphabet size, and public overlap floor jointly, with numerical comparison constants independent of all four indices over their stated domains. This uniformity makes the dependence on diminishing overlap part of the risk characterization itself.

The nearest two-channel result is Theorem 4 of CausalSmith working paper CSWP-2026-022 \citep{CausalSmith2026Annotation}. At each fixed overlap floor it gives, with constants allowed to depend on that floor,
\[
\min\left\{1,\frac1n+\frac{d^2}{(n+m)^2[\log(\exp(1)n)]^2}\right\}
\]
as the risk order for the complete-plus-auxiliary experiment. The present \leanref{S-1}{passive two-channel experiment} likewise combines complete records with independent records containing treatment and covariate information, under \cref{obj:ass:data-product}. Within the overlap class in \cref{obj:def:model}, the two budgets may be unequal: complete records supply outcome information, while both channels supply information about the treatment--covariate distribution. The overlap-uniform comparison in \cref{obj:thm:uniform-annotation-frontier}, together with the rare-label benchmark in \cref{obj:thm:rare-label-floor}, gives these information sources separate roles with numerical constants uniform in the possibly diminishing floor. The benchmark specializations are collected in \cref{obj:prop:benchmark-recovery}. Their sequence implications in \cref{obj:thm:annotation-resource-phases} characterize minimax consistency, attainment of the rare-label benchmark order, and strict relative improvement from auxiliary records as the budgets, alphabet size, and overlap floor vary together.

Reciprocal-polynomial approximation has an important antecedent in \citet[Section 3.3, Theorems 5--6]{MaZhuJiaoWainwright2022OPE}. They study off-policy evaluation with supplied target-policy weights and partial knowledge of the behavior policy through a global lower bound on action probabilities. Their Chebyshev construction and approximation-based lower bound provide methodological precedents for estimating weighted outcome functionals with uncertain denominators. The upper-bound result uses a Poisson sampling model and a delocalization condition on the target policy, while the lower bound imposes its stated exploration-probability regime. Here the target uses unknown population covariate weights, and the overlap restriction controls treatment probabilities conditional on an occupied covariate cell. Consequently, small population cells remain admissible under fixed conditional overlap. The estimator in \cref{obj:def:estimator-handle} adapts reciprocal approximation to these unknown weights and to the distinct outcome and marginal-information budgets.

The passive experiment also belongs alongside semi-supervised treatment-effect inference. \citet[Theorem 2.1 and Corollary 2.1]{ChakraborttyDai2024Semisupervised} study labeled and outcome-unlabeled samples and obtain robustness and efficiency results under their nuisance-estimation conditions. \citet[Theorems 3.7, 3.13, and 5.1]{ZhangChakraborttyBradic2023DecayingOverlap} analyze missing-at-random labeling with decaying labeling probabilities, making the observation mechanism central to the available information. A complementary design question is addressed by \citet[Theorems 4.2 and 5.1 of the arXiv manuscript]{NwankwoGoldkindZhou2026Annotations}: their annotation-allocation result optimizes asymptotic variance, and their feasible procedure uses batch adaptation and cross-fitting. Those conclusions operate under their proportional budget and batch regime, causal identification and missing-at-random conditions, treatment and annotation positivity, and the stipulated moment, nuisance-consistency, marginal and product convergence-rate, function-class, and cross-batch dependence conditions. In our experiment, independent sampling from a common population fixes how records enter each channel. The squared-error comparison therefore characterizes the value of auxiliary records under passive acquisition, with treatment overlap and the two sample sizes specified as public experiment indices.

More broadly, limited-overlap precision and inference motivate attention to the information carried by rare treatment observations \citep{Rothe2017LimitedOverlap,ArmstrongKolesar2021ATEInference,MaWang2020RobustIPW}. Discrete functional estimation supplies complementary tools for aggregate targets when many cells have small counts, particularly polynomial approximation and moment matching \citep{JiaoVenkatHanWeissman2015Functionals,WuYang2016Entropy}. These strands meet in the paper's minimax squared-error question: how treatment overlap and the separate budgets for outcome and treatment--covariate information jointly determine the precision of a population-weighted treatment contrast.
